\documentclass[../../../include/open-logic-section]{subfiles}

\begin{document}

\olfileid[id]{sth}{cardinals}{zfc}	
\olsection{$\ZFC$: Sebuah Tonggak}

Dengan penambahan Pengurutan Baik, kita telah mencapai tonggak teoretis
terakhir. Kini kita mempunyai semua aksioma yang diperlukan bagi~$\ZFC$.
Secara terperinci:

\begin{defn}
Teori $\ZFC$ mempunyai aksioma-aksioma berikut: Ekstensionalitas, Gabungan,
Pasangan, Himpunan Kuasa, Ketakhinggaan, Fondasi, Pengurutan Baik, serta semua
instans skema Pemisahan dan Penggantian. Dengan kata lain, $\ZFC$ menambahkan
Pengurutan Baik kepada~$\ZF$.
\end{defn}

$\ZFC$ merupakan singkatan dari teori himpunan
\emph{Zermelo--Fraenkel} dengan \emph{Pilihan}. Hal ini mungkin tampak agak
aneh, sebab aksioma yang kita tambahkan disebut ``Pengurutan Baik'', bukan
``Pilihan''. Namun, ketika kelak kita merumuskan {Pilihan}, ternyata
Pengurutan Baik ekuivalen (modulo~$\ZF$) dengan Pilihan (lihat
\olref[choice][woproblem]{thmwochoice}). Jadi, aksioma mana yang dipilih
sebagai aksioma ``dasar'' tidak menjadi soal. Nama ``$\ZFC$'' sendiri
sepenuhnya baku dalam literatur.

\end{document}
